% !TeX root = main.tex

% Аргументы: имя файла в images, подпись, метка для ссылки.
% После добавления файла рамка автоматически заменяется рисунком.
% Используется стандартное размещение figure, без изменения preamble.tex.
\newcommand{\imageplaceholder}[3]{%
  \begin{figure}[!htbp]
    \centering
    \IfFileExists{images/#1}{%
      \includegraphics[width=0.95\textwidth,keepaspectratio]{#1}%
    }{%
      \fbox{%
        \parbox[c][3cm][c]{0.9\textwidth}{%
          \centering
          Место для рисунка\\[0.4em]
          \texttt{images/\detokenize{#1}}
        }%
      }%
    }
    \caption{#2}
    \label{#3}
  \end{figure}
}

% Заменять вызовы команды готовым текстом по мере выполнения работы.
\newcommand{\reporttodo}[1]{%
  \par\noindent\textit{[Заполнить: #1.]}\par
}

\tableofcontents
\newpage

\section{Цель работы}

Освоить принцип адаптивной компенсации внешнего возмущения
при стабилизации многомерного линейного объекта. Исследовать влияние
коэффициента адаптации на переходные процессы состояния и управления.

\section{Постановка задачи и исходные данные}

Рассматривается объект с согласованным внешним возмущением
\begin{equation}
  \dot{\bm x}=A\bm x+\bm b u+\bm d f,
  \qquad \bm d=\bm b,
  \qquad \bm x\in\mathbb R^2.
  \label{eq:plant}
\end{equation}
Вектор состояния $\bm x(t)$ доступен измерению, матрицы объекта известны.
Возмущение $f(t)$ не измеряется; его амплитуда, частота и фаза
неизвестны регулятору.

Заданы
\begin{equation}
  A=\begin{pmatrix}3&1\\1&1\end{pmatrix},
  \qquad
  \bm b=\bm d=\begin{pmatrix}1\\2\end{pmatrix},
  \qquad t_{\text{п}}=0{,}16\,\text{с},
  \qquad \sigma=0\%.
  \label{eq:variant}
\end{equation}
Используются функция возмущения
и коэффициенты полинома фильтра:
\begin{equation}
  f(t)=\sin t,
  \qquad k_{f1}=0{,}2,
  \qquad k_{f0}=0{,}01.
  \label{eq:disturbance-data}
\end{equation}
Функция $f(t)$ используется для моделирования объекта;
в алгоритм компенсации её истинное значение не передаётся.

Требуется построить управление, обеспечивающее ограниченность
сигналов замкнутой системы и выполнение условия
\begin{equation}
  \lim_{t\to\infty}\norm{\bm x(t)}=0.
  \label{eq:goal}
\end{equation}

\section{Синтез стабилизирующей обратной связи}

\subsection{Проверка управляемости объекта}

Матрица управляемости имеет вид
\begin{equation}
  \mathcal C=\begin{pmatrix}\bm b&A\bm b\end{pmatrix}.
  \label{eq:controllability}
\end{equation}
\reporttodo{подставить матрицы варианта, вычислить определитель
и ранг $\mathcal C$, сделать вывод об управляемости}

\subsection{Выбор желаемой динамики}

Для апериодического процесса рассматривается стандартный полином
\begin{equation}
  D_{\text{ж}}(s)=(s+\omega_0)^2,
  \qquad
  A_{\text{ж}}=
  \begin{pmatrix}0&1\\-\omega_0^2&-2\omega_0\end{pmatrix},
  \qquad \omega_0=\frac{t_{\text{п}}^*}{t_{\text{п}}}.
  \label{eq:desired-dynamics}
\end{equation}
\reporttodo{указать полосу допуска для определения времени переходного
процесса, найти нормированное время $t_{\text{п}}^*$,
рассчитать $\omega_0$ и численные элементы $A_{\text{ж}}$}

\subsection{Расчёт матрицы обратной связи}

Матрица $K$ определяется методом модального управления:
\begin{equation}
  AM-MA_{\text{ж}}=\bm b H,
  \qquad K=HM^{-1},
  \qquad A_M=A-\bm b K.
  \label{eq:modal-control}
\end{equation}
Матрица $H$ выбирается так, чтобы пара $(A_{\text{ж}},H)$
была полностью наблюдаемой.

\reporttodo{выбрать $H$, привести решение уравнения Сильвестра $M$,
проверить его невырожденность, записать $K$ и $A_M$;
проверить собственные значения $A_M$}

\noindent\textbf{Вывод по пункту:}
\reporttodo{сформулировать результат проверки управляемости
и синтеза стабилизирующей обратной связи}

\section{Наблюдатель состояния модели возмущения}

\subsection{Модель генератора и вспомогательный фильтр}

Для возмущения~\eqref{eq:disturbance-data} модель генератора имеет вид
\begin{equation}
  \ddot f+f=0.
  \label{eq:disturbance-generator}
\end{equation}
Вспомогательный фильтр второго порядка задаётся матрицами
\begin{equation}
  K_f(s)=s^2+k_{f1}s+k_{f0},
  \qquad
  A_{0f}=\begin{pmatrix}0&1\\-k_{f0}&-k_{f1}\end{pmatrix},
  \qquad \bm b_{0f}=\begin{pmatrix}0\\1\end{pmatrix}.
  \label{eq:disturbance-filter}
\end{equation}
\reporttodo{подставить коэффициенты фильтра, проверить гурвицевость
$A_{0f}$ и записать вектор истинных параметров $\bm\theta_f$
для анализа результатов}

\subsection{Построение наблюдателя}

Используется наблюдатель с условием согласования
$N\bm b=\bm b_{0f}$:
\begin{equation}
  \begin{aligned}
    \hat{\bm\xi}_f&=\bm\eta+N\bm x,\\
    \dot{\bm\eta}
      &=A_{0f}\bm\eta+(A_{0f}N-NA)\bm x-N\bm b u.
  \end{aligned}
  \label{eq:observer}
\end{equation}
\reporttodo{выбрать матрицу $N$ для объекта данной работы,
проверить $N\bm b=\bm b_{0f}$ и привести численные уравнения наблюдателя}

\noindent\textbf{Вывод по пункту:}
\reporttodo{указать, за счёт какого свойства фильтра затухает
ошибка наблюдения и какие измерения использует наблюдатель}

\section{Адаптивный компенсатор возмущения}

\subsection{Закон управления и алгоритм адаптации}

Управление складывается из стабилизирующей и компенсирующей частей:
\begin{equation}
  u=-K\bm x-\hat f,
  \qquad
  \hat f=\hat{\bm\theta}_f^{\mathsf T}\hat{\bm\xi}_f.
  \label{eq:control}
\end{equation}
Матрица $P=P^{\mathsf T}>0$ находится из уравнения Ляпунова
\begin{equation}
  A_M^{\mathsf T}P+PA_M=-Q,
  \qquad Q=I_2.
  \label{eq:lyapunov}
\end{equation}
Алгоритм настройки параметров:
\begin{equation}
  \dot{\hat{\bm\theta}}_f
    =\gamma\hat{\bm\xi}_f\bm b^{\mathsf T}P\bm x,
  \qquad \gamma>0,
  \qquad \hat{\bm\theta}_f(0)=\bm 0.
  \label{eq:adaptation}
\end{equation}
\reporttodo{рассчитать $P$ и проверить, что $P>0$;
кратко пояснить свойства замкнутой системы;
при записи модели ошибок учесть начальную ошибку наблюдателя}

\clearpage
\subsection{Реализация и условия моделирования}

\reporttodo{указать начальные состояния $\bm x(0)$ и $\bm\eta(0)$,
интервал моделирования, решатель, шаг или допуски интегрирования}

Для сравнения выбираются два различных коэффициента адаптации:
\begin{equation}
  \gamma_1=\ldots,\qquad \gamma_2=\ldots,
  \qquad \gamma_1\ne\gamma_2.
  \label{eq:adaptation-gains}
\end{equation}
В обоих опытах сохраняются одинаковые начальные условия,
возмущение, матрицы обратной связи и параметры наблюдателя.
Схема моделирования представлена на рисунке~\ref{fig:scheme}.

\imageplaceholder{lab11_scheme.png}
  {Схема системы с адаптивной компенсацией возмущения в Simulink}
  {fig:scheme}

\clearpage
\section{Исследование системы при первом коэффициенте адаптации}

В первом опыте используется коэффициент $\gamma=\gamma_1$
из~\eqref{eq:adaptation-gains}.
На рисунках~\ref{fig:gamma1-state} и~\ref{fig:gamma1-control}
представлены компоненты вектора состояния и управляющее воздействие.

\imageplaceholder{lab11_gamma1_x.png}
  {Компоненты состояния $x_1(t)$ и $x_2(t)$ при $\gamma=\gamma_1$}
  {fig:gamma1-state}

\imageplaceholder{lab11_gamma1_u.png}
  {Управляющее воздействие $u(t)$ при $\gamma=\gamma_1$}
  {fig:gamma1-control}

\subsection{Анализ результатов}

\reporttodo{описать переходный процесс $x_1$ и $x_2$, оценить время
затухания и остаточные колебания, отметить максимальное управление
и характер его установившегося режима}

\noindent\textbf{Вывод по пункту:}
\reporttodo{оценить достижение цели стабилизации
при первом коэффициенте адаптации по результатам моделирования}

\clearpage
\section{Исследование системы при втором коэффициенте адаптации}

Во втором опыте используется коэффициент $\gamma=\gamma_2$
из~\eqref{eq:adaptation-gains}.
Результаты приведены на рисунках~\ref{fig:gamma2-state}
и~\ref{fig:gamma2-control}.

\imageplaceholder{lab11_gamma2_x.png}
  {Компоненты состояния $x_1(t)$ и $x_2(t)$ при $\gamma=\gamma_2$}
  {fig:gamma2-state}

\imageplaceholder{lab11_gamma2_u.png}
  {Управляющее воздействие $u(t)$ при $\gamma=\gamma_2$}
  {fig:gamma2-control}

\subsection{Анализ результатов и сравнение опытов}

\reporttodo{описать результаты второго опыта и сравнить два значения
$\gamma$ по скорости затухания состояния, величине выбросов,
остаточным колебаниям и амплитуде управления}

\noindent\textbf{Вывод по пункту:}
\reporttodo{сформулировать влияние коэффициента адаптации
и обосновать выбор предпочтительного значения для данного объекта}

\section{Выводы}

\reporttodo{кратко перечислить результаты синтеза обратной связи,
наблюдателя и адаптивного компенсатора; оценить выполнение цели
стабилизации по графикам и сформулировать влияние $\gamma$}
